% © 2026 Ing. David Jaroš — CC BY-NC-ND 4.0
%
% This work is licensed under a Creative Commons Attribution-NonCommercial-NoDerivatives
% 4.0 International License (CC BY-NC-ND 4.0).
%
% See LICENSE.md for full license text.
%
% Priority 3 — B₀=B Bridge via Seeley-DeWitt a₄ Coefficient
% Concrete mechanism: spectral action on M⁴×S¹ at self-dual point

\documentclass[12pt,a4paper]{article}
\usepackage[a4paper, margin=2.5cm]{geometry}
\usepackage{amsmath,amssymb,amsthm,booktabs,tcolorbox}
\tcbuselibrary{skins,breakable}
\usepackage{hyperref}

\newtheorem{theorem}{Theorem}[section]
\newtheorem{lemma}[theorem]{Lemma}
\newtheorem{conjecture}[theorem]{Conjecture}
\theoremstyle{remark}
\newtheorem{remark}[theorem]{Remark}
\newtheorem{openq}[theorem]{Open Question}

\newcommand{\PROVED}{\textbf{[Proved]}}
\newcommand{\STD}{\textbf{[STD]}}
\newcommand{\MC}{\textbf{[MC]}}
\newcommand{\OPEN}{\textbf{[OPEN]}}
\newcommand{\Lone}{\textbf{[L1]}}
\newcommand{\Lzero}{\textbf{[L0]}}
\newcommand{\OBS}{\textbf{[OBS]}}

\title{\textbf{$B_0 = B$ Bridge via Seeley-DeWitt $a_4$ Coefficient}\\
{\large A Concrete Mechanism for Gap G137-B Closure}}
\author{Ing.\ David Jaro\v{s}}
\date{2026-05-11}

\begin{document}
\maketitle

\begin{tcolorbox}[colback=yellow!7!white,colframe=yellow!55!black,
 title={GR-normalisation correction}]
The $a_4/B$ mechanism below remains an exploratory alpha-track proposal.
Its earlier request to ``check $a_2$'' must not cite the historical
five-dimensional coefficient note as a derivation of Newton's constant.
The corrected four-dimensional proper-time/Kaluza--Klein coefficient and its
open assumptions are given in
\texttt{canonical/gr\_closure/gap\_10d\_induced\_gravity\_endgame.tex}.
\end{tcolorbox}

\begin{abstract}
We describe a concrete new mechanism for closing Gap G137-B: the computation of
the Seeley-DeWitt coefficient $a_4$ of the spectral action functional
$S_{\mathrm{spec}}[D] = \mathrm{Tr}[f(D/\Lambda)]$ for the UBT Dirac operator $D$
on $M^4\times S^1_\psi$, evaluated at the self-dual point $R_\psi\cdot\Lambda=1$.
The conjecture is that $a_4$ evaluated at this point produces the winding
partition function $Z_{\mathrm{wind}}(\tau=i) = \vartheta_3(0|i)$ as a natural
normalisation factor.  If confirmed, the quarter-power
$Z_{\mathrm{wind}}(1)^{1/4}$ enters the effective coupling $B$ in
$V_{\mathrm{eff}}(n) = n^2 - B\,n\ln n$, closing Gap G137-B without any
fitted parameter.
The algebraic decomposition
$B = N_{\mathrm{eff}}^{3/2}\cdot 2^{1/8}\cdot Z_{\mathrm{wind}}(1)^{1/4}$
(from \texttt{rogers\_ramanujan\_c3\_connection.tex}) provides the target
formula against which the $a_4$ computation must be verified.
\end{abstract}

%──────────────────────────────────────────────────────────────────────────────
\section{Background and Motivation}
\label{sec:background}
%──────────────────────────────────────────────────────────────────────────────

\subsection{The gap}

The current alpha chain gives $n^*(B_0) \approx 65$ with $B_0 = 8\pi$ (proved,
\Lone).  The phenomenological value $B_{\mathrm{phenom}} \approx 46.3$ that
selects $n^* = 137$ must come from $S[\Theta]$ without any fitted input.

The Rogers-Ramanujan analysis
(\texttt{research\_tracks/T3\_ALPHA/rogers\_ramanujan\_c3\_connection.tex})
identifies the candidate formula:
\begin{equation}
  B = N_{\mathrm{eff}}^{3/2}\cdot 2^{1/8}\cdot Z_{\mathrm{wind}}(1)^{1/4},
  \label{eq:B_candidate}
\end{equation}
where $Z_{\mathrm{wind}}(1) = \vartheta_3(0|i)$ is the winding partition function
at the self-dual radius.  The factors $N_{\mathrm{eff}}^{3/2}$ (algebraic, \Lzero)
and $2^{1/8}$ (from $\vartheta_3/\eta = \sqrt{2}$, \STD) are already identified;
the missing piece is why $Z_{\mathrm{wind}}(1)^{1/4}$ enters $V_{\mathrm{eff}}$.

\subsection{The spectral action approach}

The UBT field equation is derived from an action $S[\Theta]$.
The spectral action formalism of Connes-Chamseddine provides a systematic way
to compute one-loop corrections to any action from the associated Dirac operator.
The Seeley-DeWitt expansion gives:
\begin{equation}
  \mathrm{Tr}\!\left[f\!\left(\frac{D}{\Lambda}\right)\right]
  \sim a_0\Lambda^4 + a_2\Lambda^2 + a_4\ln\Lambda + \cdots
  \label{eq:spectral_expansion}
\end{equation}
where $a_k$ are the Seeley-DeWitt coefficients determined entirely by the
geometry of the manifold and the Dirac operator $D$.

%──────────────────────────────────────────────────────────────────────────────
\section{The Geometry: $M^4\times S^1_\psi$}
\label{sec:geometry}
%──────────────────────────────────────────────────────────────────────────────

\subsection{Setup}

UBT defines complex time $\tau = t + i\psi$ with the imaginary component $\psi$
compactified on a circle $S^1_\psi$ of radius $R_\psi$:
\begin{equation}
  \psi \sim \psi + 2\pi R_\psi.
\end{equation}
The total manifold is $\mathcal{M} = M^4\times S^1_\psi$.

The UBT Dirac operator in this geometry is:
\begin{equation}
  D = D_{M^4} + i\frac{\partial}{\partial\psi}\otimes\gamma_5,
  \label{eq:D_ubt}
\end{equation}
where $D_{M^4}$ is the standard Dirac operator on the four-dimensional
Riemannian manifold $M^4$.

\subsection{Winding spectrum}

The eigenvalues of $i\partial_\psi$ on $S^1_\psi$ are:
\begin{equation}
  \lambda_n = \frac{n}{R_\psi}, \qquad n\in\mathbb{Z}.
\end{equation}
The trace over the $S^1_\psi$ spectrum:
\begin{equation}
  \mathrm{Tr}_{S^1_\psi}\!\left[f\!\left(\frac{\lambda_n}{\Lambda}\right)\right]
  = \sum_{n\in\mathbb{Z}} f\!\left(\frac{n}{R_\psi\Lambda}\right).
  \label{eq:trace_S1}
\end{equation}

\subsection{The self-dual point}

At $R_\psi\cdot\Lambda = 1$ (i.e.\ the T-duality self-dual point), the
Poisson resummation of~\eqref{eq:trace_S1} gives:
\begin{equation}
  \sum_{n\in\mathbb{Z}} f\!\left(\frac{n}{\Lambda R_\psi}\right)
  \xrightarrow{R_\psi\Lambda\,=\,1}
  \sum_{n\in\mathbb{Z}} f(n).
  \label{eq:poisson_self_dual}
\end{equation}
For a Gaussian test function $f(x) = e^{-\pi x^2}$, the Poisson summation
gives a theta function:
\begin{equation}
  \sum_{n\in\mathbb{Z}} e^{-\pi n^2} = \vartheta_3(0|i).
  \label{eq:theta3_from_poisson}
\end{equation}
This is the winding partition function at $\tau = i$.

%──────────────────────────────────────────────────────────────────────────────
\section{The $a_4$ Coefficient Conjecture}
\label{sec:conjecture}
%──────────────────────────────────────────────────────────────────────────────

\subsection{Standard Seeley-DeWitt $a_4$}

For a generalised Laplacian $P = -(g^{\mu\nu}\nabla_\mu\nabla_\nu + E)$
on a $d$-dimensional manifold, the Seeley-DeWitt coefficient $a_4(P)$ is:
\begin{equation}
  a_4(P) = \frac{1}{16\pi^2}\int_{\mathcal{M}}
  \mathrm{tr}\!\left[\frac{1}{6}R^2 - \frac{1}{6}R_{\mu\nu}R^{\mu\nu}
  + \frac{1}{30}R_{\mu\nu\rho\sigma}R^{\mu\nu\rho\sigma}
  + \frac{1}{2}E^2 - \frac{1}{6}RE + \frac{1}{12}\Omega_{\mu\nu}\Omega^{\mu\nu}
  \right]\sqrt{g}\,d^dx,
  \label{eq:a4_standard}
\end{equation}
where $\Omega_{\mu\nu}$ is the curvature of the connection associated to $P$.

\subsection{$a_4$ on $M^4\times S^1_\psi$ for the UBT Dirac operator}

For the product geometry $M^4\times S^1_\psi$, the $a_4$ coefficient factorises
into contributions from $M^4$ and $S^1_\psi$. The $S^1_\psi$ factor at the
self-dual point $R_\psi\Lambda=1$ contributes:
\begin{equation}
  a_4^{(S^1)}\big|_{R_\psi\Lambda=1}
  \;\propto\; \sum_{n\in\mathbb{Z}} e^{-\pi n^2/(\Lambda R_\psi)^2}
  \;\xrightarrow{R_\psi\Lambda\,=\,1}\;
  \vartheta_3(0|i).
  \label{eq:a4_S1_contribution}
\end{equation}

\begin{conjecture}[Self-dual normalisation of $a_4$]
\label{conj:a4_theta3}
The Seeley-DeWitt coefficient $a_4$ of the UBT spectral action evaluated at
the self-dual point $R_\psi\cdot\Lambda = 1$ contains a factor:
\begin{equation}
  a_4\big|_{\mathrm{self\text{-}dual}}
  = C\cdot\vartheta_3(0|i),
\end{equation}
where $C$ is determined by the $M^4$ geometry and the UBT Dirac operator
structure.  This produces $Z_{\mathrm{wind}}(1) = \vartheta_3(0|i)$ as the
natural normalisation constant in the spectral action.
\end{conjecture}

%──────────────────────────────────────────────────────────────────────────────
\section{The Bridge to $V_{\mathrm{eff}}$ and Gap G137-B}
\label{sec:bridge}
%──────────────────────────────────────────────────────────────────────────────

\subsection{Mechanism}

If Conjecture~\ref{conj:a4_theta3} holds, then the one-loop effective potential
derived from the spectral action on $M^4\times S^1_\psi$ at the self-dual point
contains a factor:
\begin{equation}
  V_{\mathrm{eff}} \supset
  N_{\mathrm{eff}}^{3/2}\cdot a_4^{(S^1)}\big|_{\mathrm{self\text{-}dual}}
  \cdot n\ln n
  = N_{\mathrm{eff}}^{3/2}\cdot\vartheta_3(0|i)^{1/4}\cdot 2^{1/8}\cdot n\ln n.
\end{equation}
This would directly produce the candidate formula~\eqref{eq:B_candidate} from
the spectral action, closing Gap G137-B.

\subsection{Three-step proof strategy}

To close Gap G137-B via this mechanism, three steps are needed:

\begin{enumerate}
  \item \textbf{Compute $a_2$ for $D_{\mathrm{UBT}}$ on $M^4\times S^1_\psi$}:
    show $a_2$ reproduces the Einstein-Hilbert action.  This is a consistency
    check confirming the UBT Dirac operator is correctly specified.
    See also \texttt{research\_tracks/T3\_ALPHA/seeley\_dewitt\_coefficients.tex}.

  \item \textbf{Compute $a_4$ at the self-dual point}:
    evaluate equation~\eqref{eq:a4_S1_contribution} explicitly and show
    $a_4|_{\mathrm{self-dual}} \propto \vartheta_3(0|i)$.

  \item \textbf{Extract $B$ from the $n\ln n$ term}:
    derive the Coleman-Weinberg effective potential from the one-loop
    contribution of the spectral action and identify the coefficient of
    $n\ln n$ as $B = N_{\mathrm{eff}}^{3/2}\cdot 2^{1/8}\cdot\vartheta_3(0|i)^{1/4}$.
\end{enumerate}

%──────────────────────────────────────────────────────────────────────────────
\section{Falsification Conditions}
\label{sec:falsification}
%──────────────────────────────────────────────────────────────────────────────

This mechanism is falsified if any of the following holds:

\begin{itemize}
  \item $a_4$ at the self-dual point does not contain a factor of $\vartheta_3(0|i)$
    (i.e.\ the product geometry $M^4\times S^1_\psi$ does not produce
    theta-function normalisation).

  \item The $n\ln n$ term in $V_{\mathrm{eff}}$ comes from $M^4$ only, with no
    $S^1_\psi$ contribution to the coefficient.

  \item The $a_4$ coefficient from the UBT Dirac operator gives a value of $B$
    inconsistent with $B_{\mathrm{phenom}}\approx 46.298$ by more than $0.1\%$.
\end{itemize}

%──────────────────────────────────────────────────────────────────────────────
\section{Summary and Next Steps}
\label{sec:summary}
%──────────────────────────────────────────────────────────────────────────────

\begin{tcolorbox}[colback=green!4!white, colframe=green!55!black,
  title={Result summary}]
\begin{center}
\begin{tabular}{lll}
\toprule
\textbf{Item} & \textbf{Status} & \textbf{Note} \\
\midrule
$Z_{\mathrm{wind}}(\tau=i) = \vartheta_3(0|i)$ & \STD & winding partition function \\
$a_4^{(S^1)}|_{\mathrm{self-dual}} \propto \vartheta_3(0|i)$ &
  \MC/\OPEN\ (Conjecture~\ref{conj:a4_theta3}) & needs explicit computation \\
$a_2$ reproduces Einstein-Hilbert & consistency check (\OPEN) &
  see \texttt{seeley\_dewitt\_coefficients.tex} \\
$a_4 \to B = N^{3/2}\cdot 2^{1/8}\cdot\vartheta_3^{1/4}$ &
  \OPEN\ (Gap G137-B) & main target \\
\bottomrule
\end{tabular}
\end{center}
\end{tcolorbox}

\paragraph{Immediate next step.}
Compute $a_4(D_{\mathrm{UBT}})$ on $M^4\times S^1_\psi$ explicitly using
the standard heat-kernel formula~\eqref{eq:a4_standard}.
The $S^1_\psi$ factor at $R_\psi\Lambda=1$ must be evaluated as a theta sum
and compared with $\vartheta_3(0|i)$.
A positive result would be the first first-principles derivation of
$B$ from $S[\Theta]$, closing Gap G137-B.

\end{document}
